%% The Elicit Author installation manual — UC-050, "Install Author from the Installation Manual".
%%
%% Build it with MANUAL=installation docs/manual/build-manual.sh, which checks the configuration
%% reference against the code (UC-051 A3, NFR-021), renders it through a TeX Live container and
%% stamps it with Author's version and build date (UC-051 BR-001). See README.md.
%%
%% It is a release artifact, attached to each tagged release and never packaged into the image
%% (C-027): the reader has no Author yet. Running the authoring tool is the author's manual
%% (UC-039); provisioning the rest of a site is the platform's installation manual. This one
%% points at both and stops.
%%
%% \opt{} is checked against Author's code by the umbrella's check-properties.sh. Settings that
%% belong to the preview Survey, not to Author, are written with \optlit{} so the check does not
%% hold Author to them.

\documentclass[11pt, a4paper]{article}

\usepackage[T1]{fontenc}
\usepackage[a4paper, margin = 25mm, bottom = 28mm]{geometry}

%% The footer's product name; the brand package is the umbrella's, shared by every Elicit manual (C-028).
\newcommand{\elicitProduct}{Elicit Author}
\usepackage{elicit-brand}
\usepackage[hidelinks]{hyperref}

\hypersetup{
  pdftitle = {Elicit Author — Installing the Authoring Tool},
  pdfauthor = {Elicit Software},
  pdfsubject = {Installing and configuring Elicit Author and its preview Survey},
  colorlinks = true,
  linkcolor = elicitInkDark,
  urlcolor = elicitInkDark,
}

\graphicspath{{./}}

\begin{document}

%% ------------------------------------------------------------------ title page
\begin{titlepage}
\thispagestyle{empty}
\centering
\vspace*{3.5cm}

\includegraphics[width = 70mm]{images/elicit-logo.png}

\vspace{1.4cm}

{\color{elicitInk}\Huge\bfseries Installing Author\par}

\vspace{0.5cm}

{\color{elicitSlate}\Large The Elicit Author installation manual\par}

\vspace{0.4cm}

{\color{elicitSlate}\large Author \textperiodcentered\ preview Survey\par}

\vspace{2.2cm}

{\color{elicitSlate}\large
Version \textbf{\elicitVersion}\\[0.35em]
Built \textbf{\elicitBuildDate}\par}

\vfill

{\color{elicitSlate}\small
This manual describes the Author version named above. It is published beside that\\
release and is read before any Author service is running.\\[0.6em]
Elicit Software \textperiodcentered\ PolyForm Noncommercial License 1.0.0\par}

\vspace{1cm}
\end{titlepage}

%% ------------------------------------------------------------------ contents
\pagenumbering{roman}
\setcounter{page}{1}
{\color{elicitInk}\section*{Contents}}
\vspace{-0.6em}
\begingroup
\hypersetup{linkcolor = elicitCharcoal}
\makeatletter\@starttoc{toc}\makeatother
\endgroup

\vfill

\begin{elicitnote}[Who this manual is for]
Someone comfortable with PostgreSQL, a container runtime and an
OIDC provider, who is adding the authoring tool to an Elicit installation or setting it up on
its own. It covers \emph{provisioning} Author. Building a survey is the author's manual, and
provisioning a site's Survey and Admin is the platform's installation manual.
\end{elicitnote}

\clearpage
\pagenumbering{arabic}
\setcounter{page}{1}

%% ================================================================== 1
\section{What Author is}
\label{sec:what}
\markboth{What Author is}{}

\term{Author} is the tool in which surveys are designed. An author lays out steps, sections and
questions, connects them with rules, and \emph{exports} the result as a survey definition
file (\file{.elicit}). An administrator applies that file at a site through Admin's
\ui{Apply Survey Definition}. The file is the only thing that ever travels from Author to a
site: Author calls neither Survey, Admin, FHHS nor Pedigree at any site, and no site calls
Author.

\begin{elicitnote}[A site that only delivers surveys needs no Author]
A site that applies definition files written elsewhere runs Survey and Admin and nothing
from this manual. The platform's installation manual covers that site completely. Install
Author where surveys are written, which may be one place serving many sites.
\end{elicitnote}

\subsection{What Author needs}

\begin{plaintable}{tab:needs}{What an Author installation is made of}{@{}>{\raggedright\arraybackslash}p{0.22\linewidth} >{\raggedright\arraybackslash}p{0.72\linewidth}@{}}{\thead{Piece} & \thead{What it is}}
\ui{Author} & The authoring application, image \file{elicitsoftware/author}. Signed-in
authors work in it from a browser. \\
\addlinespace
A database & Author's own database, called \val{author} here, carrying the same
\val{survey} schema a site's database carries. Author writes working copies of surveys into
it. It is never a site's survey database (Chapter~\ref{sec:database}). \\
\addlinespace
The preview Survey & A second instance of the Survey application, image
\file{elicitsoftware/survey}, pointed at the \val{author} database. It builds the
\val{survey} schema there on its first start, and authors walk through a survey as a
respondent would (Chapter~\ref{sec:preview}). \\
\addlinespace
An OIDC client & Author's client at the identity provider, \val{elicit-author}, with two
roles (Chapter~\ref{sec:oidc}). \\
\addlinespace
Brand and translations & The same brand directory and curated translations the site
applications mount (Chapter~\ref{sec:brand}). Optional: without them Author runs in the
default brand and in English. \\
\end{plaintable}

\subsection{The start order}

The database first, then the preview Survey, then Author. The preview Survey creates the
schema Author maps its entities onto. Author never creates it and does not check it at start,
so an Author started against an empty \val{author} database comes up, but every screen fails
and the \ui{System} database screen reports the schema missing until the preview has started
once.

%% ================================================================== 2
\section{Before you start}
\label{sec:before}
\markboth{Before you start}{}

\begin{itemize}
  \item \textbf{PostgreSQL 17}, with the Elicit database roles (\val{dbowner}, \val{dbuser},
        \val{elicit_owner}, \val{survey_user}). A site's cluster already has them; the
        platform's installation manual gives the role script for a cluster that does not.
  \item \textbf{An OIDC provider} on which you can create a confidential client and two
        client roles. Keycloak is what Elicit is tested with.
  \item \textbf{A container runtime}. Author and the preview Survey are released as container
        images only; there is no separate Java distribution to install.
  \item \textbf{The Survey image of the same release} as Author, for the preview. Author and
        the preview are tested together at one version.
  \item The addresses a browser will reach Author and the preview at, for the OIDC redirect
        and for the preview links.
\end{itemize}

%% ================================================================== 3
\section{The database}
\label{sec:database}
\markboth{The database}{}

Author works in a database of its own. It may live in the same cluster as a site's
\val{survey} database, but it is never that database: authors save half-finished work, and
nothing they do should touch respondent data.

Run this once, as a superuser, in the cluster. It creates the database, grants it to the
Elicit roles, and creates the two empty schemas. Every table inside them comes from the
preview Survey's migrations on its first start.

\begin{elicitsql}
CREATE DATABASE author;
GRANT CONNECT ON DATABASE author TO dbowner;
GRANT CREATE ON DATABASE author TO dbowner;
GRANT CONNECT ON DATABASE author TO dbuser;
\c author

GRANT USAGE, CREATE ON SCHEMA public TO dbowner;
GRANT USAGE ON SCHEMA public TO dbuser;

CREATE SCHEMA survey AUTHORIZATION elicit_owner;
GRANT USAGE ON SCHEMA survey TO survey_user;
GRANT USAGE ON SCHEMA survey TO surveyadmin_user;
GRANT USAGE ON SCHEMA survey TO surveyreport_user;

CREATE SCHEMA surveyreport AUTHORIZATION elicit_owner;
GRANT USAGE ON SCHEMA surveyreport TO surveyreport_user;
GRANT USAGE ON SCHEMA surveyreport TO surveyadmin_user;
GRANT USAGE ON SCHEMA surveyreport TO survey_user;

GRANT USAGE, CREATE ON SCHEMA surveyreport TO dbowner;
GRANT SELECT, INSERT, UPDATE, DELETE
  ON ALL TABLES IN SCHEMA surveyreport TO dbowner;
ALTER DEFAULT PRIVILEGES IN SCHEMA surveyreport
  GRANT SELECT, INSERT, UPDATE, DELETE ON TABLES TO dbowner;
GRANT USAGE ON ALL SEQUENCES IN SCHEMA surveyreport TO dbowner;
ALTER DEFAULT PRIVILEGES IN SCHEMA surveyreport
  GRANT USAGE ON SEQUENCES TO dbowner;
\end{elicitsql}

Author connects as \val{survey_user}, the role the preview's migrations grant full access to
every definition table. Author owns no migrations of its own and never creates or alters the
schema.

\begin{elicitnote}[The umbrella's compose file does this for you, once]
The umbrella mounts this script as \file{db-init/V0.0.4__CREATE_AUTHOR_DATABASE.sql} into
the database container's initialization directory. That directory runs \emph{only} against an
empty data directory, so a cluster that was first started before the script was there has no
\val{author} database. Run the script once by hand against it:
\end{elicitnote}

\begin{elicitcmd}
docker compose exec -T db psql -U survey -d survey \
  < db-init/V0.0.4__CREATE_AUTHOR_DATABASE.sql
\end{elicitcmd}

%% ================================================================== 4
\section{The identity provider}
\label{sec:oidc}
\markboth{The identity provider}{}

Author authenticates against a confidential OIDC client of its own. It does not share Admin's
client, and a role on Admin's client grants nothing in Author.

\subsection{The client}

\begin{plaintable}{tab:client}{Author's OIDC client}{@{}>{\raggedright\arraybackslash}p{0.22\linewidth} >{\raggedright\arraybackslash}p{0.20\linewidth} >{\raggedright\arraybackslash}p{0.52\linewidth}@{}}{\thead{Client} & \thead{Type} & \thead{Redirect}}
\val{elicit-author} & confidential & The address a browser reaches Author at, for example
\val{http://localhost:8084/*}. Post-logout redirects return to the same address. \\
\end{plaintable}

The client needs a secret, which reaches Author as \opt{OIDC_CLIENT_SECRET}. Author accepts
browser sign-ins and bearer tokens on the same client. It signs in through
\val{/callback} and signs out through \val{/logout}, both under its own address, which the
redirect pattern above covers.

\subsection{The roles}

\begin{plaintable}{tab:roles}{The roles Author checks, both on the elicit-author client}{@{}>{\raggedright\arraybackslash}p{0.24\linewidth} >{\raggedright\arraybackslash}p{0.70\linewidth}@{}}{\thead{Role} & \thead{What it permits}}
\val{elicit_author} & Authoring: every survey screen, import and export, the author's manual
and the AI drafting page. Without it a signed-in user sees only a page saying access was not
granted. \\
\addlinespace
\val{elicit_author_admin} & The \ui{System} section: setup and diagnostics. It grants no
authoring; give a person both roles if they need both. \\
\end{plaintable}

Neither role implies the other, and neither is implied by Admin's \val{elicit_admin}. Being an
administrator of a site grants nothing in Author.

\subsection{The bundled example}

The umbrella's Keycloak realm, \file{keycloak/elicit-realm.json}, carries this client, both
roles and two accounts: \val{author}/\val{author} with \val{elicit_author}, and
\val{admin}/\val{admin} with both Author roles besides Admin's. It is an evaluation provider.

\begin{elicitsecurity}
The bundled client's secret is the literal string \val{secret} and the accounts' passwords
are their names. Create the client and accounts in your own provider, with a generated
secret, before Author reaches anyone.
\end{elicitsecurity}

%% ================================================================== 5
\section{The preview Survey}
\label{sec:preview}
\markboth{The preview Survey}{}

The preview is an ordinary Survey application of the same release, configured for the
\val{author} database. It does two jobs:

\begin{itemize}
  \item On its first start its migrations build the \val{survey} schema in the \val{author}
        database. Author has no migrations of its own.
  \item Authors use it to walk through a survey as a respondent would. Author also calls it
        to drop a survey's reporting schema before deleting the survey; if it cannot be
        reached, that delete is refused with the reason shown.
\end{itemize}

It needs Survey's usual configuration (the platform's installation manual lists it), with
these differences:

\begin{plaintable}{tab:preview}{How the preview differs from a site's Survey}{@{}>{\raggedright\arraybackslash}p{0.40\linewidth} >{\raggedright\arraybackslash}p{0.54\linewidth}@{}}{\thead{Survey setting} & \thead{For the preview}}
\optlit{quarkus.datasource.jdbc.url} & The \val{author} database, not a site's
\val{survey} database. \\
\addlinespace
\optlit{quarkus.datasource.owner.jdbc.url} & The same: the migrating connection must build
the schema in \val{author}. \\
\addlinespace
\optlit{accessCode.autoRegister} & \val{true}: an author types any access code to walk
through a draft. Never set this on a site's Survey. \\
\addlinespace
\optlit{elicit.etl.enabled} & \val{false}: nothing reports on preview answers, so no
reporting schema is built for them. \\
\addlinespace
\optlit{elicit.etl.drop.enabled} & \val{true}: lets Author drop a survey's reporting schema
before deleting it, should one exist. \\
\addlinespace
\optlit{ELICIT\_SERVICE\_TOKEN} & The same value Author is given
(Section~\ref{sec:servicetoken}). \\
\end{plaintable}

The preview connects as \val{survey_user} and migrates as \val{elicit_owner}, with the same
passwords those roles have in the cluster (\optlit{ELICIT\_DB\_PASSWORD} and
\optlit{ELICIT\_OWNER\_DB\_PASSWORD}). Author reaches it at \opt{author.preview.survey.url}, an
address on the network Author runs on; it need not be reachable from a browser for that call.

\subsection{The service token}
\label{sec:servicetoken}

Elicit's applications authenticate their calls to each other with one shared secret per
deployment, the service token. Author makes one such call, the request that drops a reporting
schema, and sends the token with it as a bearer credential; the preview Survey answers that
request only when it carries the token the preview was given itself. Both containers therefore
receive the same environment variable, \opt{ELICIT_SERVICE_TOKEN}. An Author installed apart
from a site has a token of its own, not the site's.

Generate it once, and supply it the way the database passwords are supplied (the platform's
secret store, a container secret, or the \file{.env} file beside the compose file), never in
a file under version control:

\begin{elicitcmd}
openssl rand -hex 32
\end{elicitcmd}

\begin{itemize}
  \item The token is at least 32 characters. An application given a shorter one does not
        start, and its log names the variable and the length required.
  \item An application given no token starts, logs one warning, and reports not ready:
        \val{/q/health/ready} carries a check named \val{service-token} with the reason.
        Author then refuses to delete a survey that has a reporting schema rather than
        make the call without a credential, and the preview answers the request with
        \val{503}.
  \item Author sends the token to the address in \opt{author.preview.survey.url} and nowhere
        else, never follows a redirect with it, and never writes it to a log, a page or an
        error message. The \ui{System} overview reports it as present or absent only.
  \item To replace it, give both containers the new value and restart them together.
\end{itemize}

\begin{elicitsecurity}
The preview accepts any access code and is meant for authors' browsers only. Publish it to
authors alone, behind the same network boundary or sign-on as Author. The token protects the
request Author makes to it, not the preview's pages.
\end{elicitsecurity}

%% ================================================================== 6
\section{Installing with Docker Compose}
\label{sec:compose}
\markboth{Installing with Docker Compose}{}

The Elicit umbrella repository runs the whole platform from one compose file, Author
included: \val{author} on port 8084 and \val{author-survey}, the preview, on 8085. On a fresh
data directory it creates the \val{author} database (Chapter~\ref{sec:database}), and its
realm already carries Author's client and roles.

\begin{elicitcmd}
git clone https://github.com/ElicitSoftware/Elicit.git
cd Elicit
./cloneAllProjects.sh
docker compose up -d
\end{elicitcmd}

Authoring needs none of the site applications. To run only Author's side:

\begin{elicitcmd}
docker compose up -d author keycloak jaeger
\end{elicitcmd}

\val{author} waits for the database and for the preview to report ready, so the start order
of Chapter~\ref{sec:what} is kept for you. Every application image is pulled at one tag,
\opt{ELICIT_TAG} (default \val{latest}), so Author and its preview are always the same
release:

\begin{elicitcmd}
ELICIT_TAG=3.0.0-beta.1 docker compose up -d author keycloak jaeger
\end{elicitcmd}

\begin{elicitsecurity}
The compose file's database passwords (\val{SURVEYPW}), client secret (\val{secret}) and
service token (\val{elicit-dev-service-token-not-for-a-real-site}, used when
\opt{ELICIT_SERVICE_TOKEN} is unset) are development values. Author logs a warning at startup
and shows one on the \ui{System} overview while that token is in use. Replace them, and the
bundled realm, before Author reaches anyone
(Chapter~\ref{sec:oidc}).
\end{elicitsecurity}

%% ================================================================== 7
\section{Installing without Compose}
\label{sec:images}
\markboth{Installing without Compose}{}

Author and the preview are released as container images only, \file{elicitsoftware/author}
and \file{elicitsoftware/survey}, at the same tag. A site that runs its containers under
something other than the umbrella's compose file (Podman, Kubernetes, a container service)
starts the same two images against its own database and identity provider. Configuration
reaches them as environment variables (Chapter~\ref{sec:config}); nothing ever orders their
start for you.

\begin{enumerate}
  \item Provision the \val{author} database (Chapter~\ref{sec:database}).
  \item Start the preview Survey with the settings of Chapter~\ref{sec:preview}, the
        service token among them, and wait
        until \val{/q/health/ready} answers \val{200}.
  \item Start Author, mounting the brand and the translations read-only:
\end{enumerate}

\begin{elicitcmd}[fontsize = \footnotesize]
JDBC_URL='jdbc:postgresql://db.example.org:5432/author'\
'?connectTimeout=30&socketTimeout=600&tcpKeepAlive=true&tcpKeepAlivesIdle=60'
docker run -d --name author -p 8084:8080 \
  -e QUARKUS_DATASOURCE_JDBC_URL="$JDBC_URL" \
  -e ELICIT_DB_PASSWORD='<survey_user password>' \
  -e QUARKUS_OIDC_AUTH_SERVER_URL='https://idp.example.org/realms/elicit/' \
  -e OIDC_CLIENT_SECRET='<elicit-author client secret>' \
  -e ELICIT_SERVICE_TOKEN='<the token the preview was given>' \
  -e AUTHOR_PREVIEW_SURVEY_URL='http://preview.internal:8080' \
  -e BRAND_FILE_SYSTEM_PATH=/opt/brand \
  -e I18N_FILE_SYSTEM_PATH=/opt/i18n \
  -v /srv/elicit/brand:/opt/brand:ro \
  -v /srv/elicit/i18n:/opt/i18n:ro \
  elicitsoftware/author:<version>
\end{elicitcmd}

\val{<version>} is the version on this manual's title page; the preview Survey runs at the
same tag. Author listens on \val{8080} inside its container unless \val{QUARKUS_HTTP_PORT} says otherwise. Behind a
reverse proxy it trusts the forwarded host, prefix and protocol headers, so the OIDC
redirect is built from the address the browser used.

%% ================================================================== 8
\section{Configuration}
\label{sec:config}
\markboth{Configuration}{}

Every setting can be given as a system property, or as an environment variable named by
upper-casing it and replacing each dot and hyphen with an underscore:
\opt{author.preview.survey.url} becomes \val{AUTHOR_PREVIEW_SURVEY_URL}. A few settings read
a shorter variable of their own, named in the table. The \emph{Change} column says whether a
site must set the value (\chgset), may (\chgmay), or should leave it alone (\chgkeep).

\subsection{Database and identity provider}

\begin{settings}{tab:cfg-core}{Database and identity-provider settings}
\opt{quarkus.datasource.jdbc.url} & \val{jdbc:postgresql://db:5432/survey} & \chgset &
The database Author works in. The shipped default names a database called \val{survey};
point it at \val{author}, keeping the default's connection options, which
Chapter~\ref{sec:images} shows in full. \\
\opt{quarkus.datasource.username} & \val{survey_user} & \chgkeep & The role Author connects
as. \\
\opt{quarkus.datasource.password} & --- & \chgset & Its password. Environment variable
\opt{ELICIT_DB_PASSWORD}. A secret. \\
\opt{quarkus.oidc.auth-server-url} & \val{http://host.docker.internal:8180/realms/elicit/} &
\chgset & The identity provider's realm. It must resolve both from Author and from the
browser, because the browser is redirected to it. \\
\opt{quarkus.oidc.client-id} & \val{elicit-author} & \chgmay & Author's client. \\
\opt{quarkus.oidc.credentials.secret} & --- & \chgset & The client's secret. Environment
variable \opt{OIDC_CLIENT_SECRET}. A secret. \\
\end{settings}

\subsection{Authoring}

\begin{settings}{tab:cfg-author}{Authoring settings}
\opt{author.preview.survey.url} & \val{http://localhost:8085} & \chgset & The preview
Survey, at an address Author can reach (Chapter~\ref{sec:preview}). \\
\opt{elicit.service.token} & --- & \chgset & The service token Author sends to the preview
Survey, the same value the preview is given (Section~\ref{sec:servicetoken}). At least 32
characters. Environment variable \opt{ELICIT_SERVICE_TOKEN}. A secret. \\
\opt{author.content.languages} & \val{en,ar,es-419} & \chgkeep & The languages a survey's
content may be published in. Set by the release to the curated translations and checked
against them when the release is built; see Chapter~\ref{sec:brand}. \\
\opt{author.manual.resource} & \val{manual/elicit-author-manual.pdf} & \chgkeep & Where the
author's manual sits on the classpath. \\
\opt{author.ai.specification.resource} & \val{ai/AUTHORING_FROM_AI.md} & \chgkeep & Where
the AI authoring specification sits on the classpath. \\
\end{settings}

\subsection{Brand and translations}

\begin{settings}{tab:cfg-brand}{Brand and translation settings}
\opt{brand.file.system.path} & \val{/brand} & \chgmay & The mounted brand directory.
Missing assets fall back to the default brand. \\
\opt{brand.local.path} & \val{brand} & \chgkeep & A second, relative brand directory
consulted before the default brand. \\
\opt{i18n.file.system.path} & \val{/i18n} & \chgmay & The mounted translations; Author reads
its own \val{author/} sub-directory of it. \\
\opt{i18n.app.name} & \val{author} & \chgkeep & That sub-directory's name. \\
\opt{i18n.offered.locales} & --- & \chgmay & Offer fewer languages than are mounted, for
example \val{en,es-419}. Unset offers every mounted language. \\
\opt{i18n.pseudo-locale.enabled} & \val{false} & \chgkeep & A test-only language that shows
every text's key. \\
\end{settings}

\subsection{Operation}

\begin{settings}{tab:cfg-ops}{Operational settings}
\opt{quarkus.log.level} & \val{WARN} & \chgmay & The root log level. Environment variable
\opt{LOG_LEVEL}; Author's own classes follow \opt{ELICIT_LOG_LEVEL} (default \val{INFO}). \\
\opt{quarkus.http.cors.origins} & --- & \chgmay & Origins allowed to call Author from another
site's pages. Environment variable \opt{CORS_ALLOWED_ORIGINS}. None by default, and none is
needed: Author's pages call only Author. \\
\opt{quarkus.application.name} & \val{elicit-author} & \chgmay & The service name in metrics
and traces. \\
\opt{build.timestamp} & build & \chgkeep & Stamped by the build and shown on the System
overview. \\
\opt{elicit.debug.reveal-tokens} & \val{false} & \chgkeep & Shows raw tokens on the OIDC
diagnostics screen. Never set it outside a development machine. \\
\end{settings}

Telemetry follows Quarkus's OpenTelemetry settings, \val{quarkus.otel.*}. Every setting
this manual does not list is in the Quarkus configuration reference:
\url{https://quarkus.io/guides/all-config}.

%% ================================================================== 9
\section{Brand and translations}
\label{sec:brand}
\markboth{Brand and translations}{}

Author mounts the \emph{same} brand directory as the site applications, so the authoring tool
looks like the rest of the installation. A mounted brand may be partial; whatever it lacks
comes from the default brand.

The image carries English only. Every other language is a file in the curated translations
ElicitSoftware publishes beside each release, the umbrella's \file{elicit-i18n/}, mounted
read-only. Author reads its own \val{author/} sub-directory of it, as Survey reads
\val{survey/} and Admin \val{admin/}. The platform's installation manual explains both
mechanisms, fallback and right-to-left layout included.

Two settings are easily confused:

\begin{description}
  \item[\opt{i18n.offered.locales}] Which languages \emph{Author's own screens} are offered
        in at this installation. A site may narrow it.
  \item[\opt{author.content.languages}] Which languages \emph{a survey's content} may be
        published in. The release sets it to the curated set, and the release's build fails
        if it differs, because content in a language a site's Survey lacks could never reach a
        respondent. Leave it.
\end{description}

%% ================================================================== 10
\section{Verifying the installation}
\label{sec:verify}
\markboth{Verifying the installation}{}

In this order; each check assumes the ones before it passed.

\begin{enumerate}
  \item The preview's \val{/q/health/ready} answers \val{200}.
  \item Author's \val{/q/health/ready} answers \val{200}. Both endpoints are open without
        signing in.
\end{enumerate}

\begin{elicitcmd}
curl -fsS http://localhost:8085/q/health/ready
curl -fsS http://localhost:8084/q/health/ready
\end{elicitcmd}

\begin{enumerate}
  \setcounter{enumi}{2}
  \item A user holding \val{elicit_author} signs in at Author's address and reaches the
        \ui{Surveys} list.
  \item A user holding \val{elicit_author_admin} opens \ui{System}. The overview reports every
        required setting as present and the database check passes; the database screen
        reports the \val{survey} schema and the preview's migration level.
  \item An author creates a survey, adds a question, and exports it. An administrator applies
        the file at a site through Admin's \ui{Apply Survey Definition}.
\end{enumerate}

When all five pass, hand Author to its authors with the author's manual, which Author serves
from its own \ui{Manual} link.

%% ================================================================== 11
\section{When it does not work}
\label{sec:trouble}
\markboth{When it does not work}{}

\begin{plaintable}{tab:trouble}{Failures actually seen, and their cures}{@{}>{\raggedright\arraybackslash}p{0.36\linewidth} >{\raggedright\arraybackslash}p{0.58\linewidth}@{}}{\thead{Symptom} & \thead{Cause and cure}}
Author does not start, naming \opt{ELICIT_DB_PASSWORD} or \opt{OIDC_CLIENT_SECRET}. & A
required secret is unset. Supply it; the System overview lists every required setting. \\
\addlinespace
Author starts, but every screen fails and the database screen reports the \val{survey}
schema missing. & The preview Survey has not built the schema yet, or Author and the preview
use different databases. Start the preview against the same database first. \\
\addlinespace
The database \val{author} does not exist. & The cluster was initialized before the database
script was mounted. Run it once by hand (Chapter~\ref{sec:database}). \\
\addlinespace
Sign-in works, then a page says access was not granted. & The user lacks
\val{elicit_author}, or holds it on another client. The OIDC diagnostics screen under
\ui{System} shows the roles the provider actually sent. \\
\addlinespace
Sign-in loops, or the provider rejects the redirect. & The redirect URI on
\val{elicit-author} does not cover the address the browser used, or
\opt{quarkus.oidc.auth-server-url} is not reachable from the browser. \\
\addlinespace
Deleting a survey is refused, naming a connection failure. & Author cannot reach the preview
at \opt{author.preview.survey.url}. Check the address from Author's host. \\
\addlinespace
Author or the preview is not ready, and \val{/q/health/ready} names \val{service-token}. &
\opt{ELICIT_SERVICE_TOKEN} is unset for that container. Supply it
(Section~\ref{sec:servicetoken}). \\
\addlinespace
Author or the preview does not start, naming \opt{ELICIT_SERVICE_TOKEN}. & The token is
shorter than 32 characters. Generate a longer one. \\
\addlinespace
Deleting a survey is refused, naming the service token. & Author has no token, or Author and
the preview were given different ones. Give both containers the same value and restart
them. \\
\addlinespace
A language is missing from Author's language selector. & The mounted translations lack
\val{author/translations_<tag>.properties}, or \opt{i18n.offered.locales} leaves the
language out. \\
\end{plaintable}

%% ================================================================== 12
\section{Where to look next}
\label{sec:next}
\markboth{Where to look next}{}

\begin{description}
  \item[The author's manual] Building a survey, from an empty list to an exported definition.
        Author serves it from its \ui{Manual} link.
  \item[The platform's installation manual] Provisioning a site's Survey and Admin: the
        database roles, the site's identity-provider client, branding and translations in
        depth, applying a definition, and the report services a survey's \ui{Reports\ldots}
        page lists, with FHHS and Pedigree as the worked example.
  \item[\url{https://github.com/ElicitSoftware/Author}] Author's source, its use cases under
        \file{docs/use_cases/}, and this manual's sources under \file{docs/manual/}.
\end{description}

\vspace{1em}

\begin{elicitnote}[Which version is this?]
This manual describes Elicit Author \textbf{\elicitVersion}, built
\textbf{\elicitBuildDate}. The same pair is printed at the foot of every page, so a printed
copy can always be matched back to the release it came from.
\end{elicitnote}

\end{document}
